\begin{center}
{\LARGE\bfseries\color{ink} T-Adaptive Segment Routing}\\[0.4em]
{\color{ink}\rule{3.2cm}{0.7pt}}\\[0.45em]
{\large\color{ink} A scalar maximum-link-utilization model on Set~A}\\[0.3em]
{\small\color{muted}
Vuong Minh \enspace\textbar\enspace October 2026 \enspace\textbar\enspace
Orange / EURO-ROADEF 2026}
\end{center}

\begin{abstract}
\noindent
\small
This report studies the course Orange project derived from the EURO/ROADEF 2026
$T$-Adaptive Segment Routing challenge~\citep{orange2026}.
We implement and solve the scalar problem
$\min L=\max_{t,a}\lambda(a,t)$, that is, the minimization of maximum link utilization (MLU).
MLU is the first component of the challenge's lexicographic objective and a standard
traffic-engineering criterion used to suppress hot spots and preserve residual capacity.
All Set~A instances are solved for period~0 (1800\,s) and for periods~0 and~1 with the
reconfiguration budget (3600\,s), with averages over ten seeded runs.
Bounds, gaps and the label ``MLU-optimal'' refer only to this scalar problem,
not to full lexicographic optimality.
\end{abstract}
